\subsection*{重复相邻与重复交错之间差了一步}

第四章已经给出幂等等式 $f(f(s))=f(s)$，以及交换等式 $f(g(s))=g(f(s))$。现在用它们回答更具体的问题：已经执行 f，中间执行了 g，旧 f 又来一次，怎样使最后一次重复没有影响？

在当前允许的状态集合中，假设 f 幂等，而且 f 与 g 交换。把先 f、再 g、再 f 写成 $f(g(f(s)))$，就可以先交换最外两步，再合并相邻重复：
\[
  f(g(f(s)))=g(f(f(s)))=g(f(s)).
\]
右边表示先 f、再 g，恰好是去掉迟到重复以后的历史。固定赋值与减 1000 不交换，所以前节的赋值历史无法完成第一步。加入 gift 与加入 priority 交换，两次加入 gift 又幂等，因此标签历史中可以消去一次重复。

\CThirteenFigure{ch13-interleaving-proof}{先用交换律把两次 f 移到一起，再用幂等性消去一次重复}{fig:ch13-interleaving-proof}

如果一组状态变换两两交换，每个变换又幂等，那么有限历史中的重复可以通过交换逐步靠拢，再合成一次。这里比较的是选定状态的最终值。若每次加入 gift 都另发一封邮件，邮件也进入观察，第二封邮件就会使整体变换失去幂等性。原结果若包含首次修改时间或随机编号，重试也要从相应记录取得这些字段。

两个减法在没有下限的整数余额上交换。从共同初态 10000，先减 3000 再减 1000，与反过来执行，都得到 6000。两次减 3000 则得到 4000。不同增量允许怎样重排，同一增量送达几次，都能在这几条余额变化中分别看出来。

\subsection*{合并次序与分组为什么是两项性质}

现在让 A、B、C 三个副本分别保存标签集合 $\{\mathrm{gift}\}$、$\{\mathrm{priority}\}$、$\{\mathrm{fragile}\}$。A 先与 B 合并，再与 C 合并，得到三个标签。也可以先合并 B、C，再把结果交给 A。两种分组对应结合律：
\[
  (x\cup y)\cup z=x\cup(y\cup z).
\]
结合律允许同步系统分批、分层汇总。交换律允许调换两个输入，幂等性允许同一状态重复到达。三项性质分别处理分组、顺序与重复，集合并恰好同时具备它们。

函数复合本身也满足结合律，但把两个扣款函数组成一个批次，照样会执行其中每笔扣款。同步系统采用集合并、逐分量最大值还是函数复合，决定了这里检验的究竟是哪项运算。

把上述集合状态重复传播，无论经过几条路径，同一 gift 都只有一个成员资格，接收端可以直接合并状态。若业务还要求为每次标签设置返回独立审计编号，审计对象就保存这些编号与原设置的关系，标签集合负责恢复成员资格。

\subsection*{信息增加与界面数字增加怎样区分}

为了说明一份状态是否包含另一份状态的信息，引入偏序关系，记为 $x\preceq y$。它要求状态与自己可比，两个方向都包含时视为同一状态，而且包含关系可以连续传递。偏序允许一部分状态互不可比。$\{\mathrm{gift}\}$ 与 $\{\mathrm{priority}\}$ 就互不包含，二者都被 $\{\mathrm{gift},\mathrm{priority}\}$ 包含。

对只添加集合，$\preceq$ 可以直接取集合包含关系 $\subseteq$。一次添加满足 $L\subseteq L\cup\{\mathrm{gift}\}$，即更新后的信息至少包含原来全部信息，这样的更新称为膨胀更新。若一个函数还满足 $x\preceq y$ 时 $f(x)\preceq f(y)$，则称它相对于该偏序单调。膨胀比较同一输入更新前后，单调比较两个有序输入经过同一函数以后。

下一节的状态型合并要求本地更新膨胀。副本由空集合变成 $\{\mathrm{gift}\}$，再变成 $\{\mathrm{gift},\mathrm{priority}\}$，每一步保留已有信息。合并把两个互不可比状态送到共同上方，积累双方所知。

\CThirteenFigure{ch13-original03}{集合包含形成信息偏序，两个不可比状态通过并集到达共同上方}{fig:ch13-information-order}

信息偏序与屏幕上的数字大小可以不同。先保存已确认收款 3000 这项事实，再增加已确认退款 1000 这项事实，已知事实集合扩大，净收款却从 3000 变为 2000。界面数值由多项事实计算，减少也能来自信息增加。后面用两个只增长分量表达正负变化，就沿着这条关系构造状态。

订单上的取消与发货也不能凭显示词排成一个任意大小序列。若把取消编号为 4、已发货编号为 2，取最大值会显示取消，却没有说明包裹如何停止。业务根据取消与发货的授权、先后和外部证据，决定撤销、拦截或后续处置，投影再显示这些决定及其进度。

\subsection*{收敛以前和以后都要检查业务不变量}

设 A、B 各有余额 3000 的副本，两者共同来自同一初态，尚无扣款。网络暂时不通，A 与 B 都接受一项不同的 2000 分购买，各自检查本地 3000 足够，扣后各显示 1000。最后汇总两项已接受增量，总支出为 4000，资金核对要求余额为 $-1000$。

这两项减法可以交换，最后也可以让所有副本收敛到同一负数。余额非负要求允许状态始终留在安全范围内，因而还要检查两个单独安全的分支在合并以后是否安全。

\CThirteenFigure{ch13-original04}{同一初始 3000 被两个副本分别用于 2000 分购买，汇总后余额为负}{fig:ch13-negative-balance}

若用取最大值合并两个总支出 2000，最终支出只剩 2000，负数似乎消失了，两份实际购买中却有一份没有计入账目。若每个来源各保存自己的累计支出，合并可以完整保留 4000，负余额随之显现。资金核对把每笔已接受的合法购买都算进去，才能发现这两条分支共同花掉了过多金额。

第~\ref{section:13-5}~节将为各副本预分配可扣减额度，让两个分支消耗不同份额。副本拿同一余额各自判断，会反复使用同一份购买能力。专属份额则把可独立消费的范围先分开，合并时再核对分配、转移与总支出。
\FloatBarrier
